\section{Examples of Data Synthesis Prompts}
\label{sec:data_synthesis_prompts}

\begin{promptblock}{Image Question Answering Prompt}
\begin{verbatim}
You are given an image. Your job is to create ONE high-quality multimodal training example for an Image Question Answering (IQA) dataset.
The final answer MUST be a single JSON object and nothing else.

STEP 1 - Visual Description (less or equal than 500 {language} words)
- General scene summary and object-level details (attributes, positions, relations).
- Contextual features (environment, lighting, actions).
- Brainstorm the types of reasoning enabled (e.g., spatial, comparative, predictive).

STEP 2 - Task Selection
Choose ONE task type from the list below that best fits the image content:
- Factoid Identification: Questions about specific entities, brands, or basic facts (e.g., "What brand is the watch?").
- Visual Reasoning: Questions requiring logical inference or analysis (e.g., "How many rats were fed the control diet?").
- OCR-based Data Extraction: Questions targeting text, tables, or document info (e.g., "Who is the author of the book?").
- Domain-specific Knowledge Inquiry: Questions requiring specialized background knowledge (e.g., "What style of architecture is this?").

STEP 3 - Populate the Example
Fill every key below using double quotes. Do not add extra keys.
{
 "description": "<STEP 1 output>",
 "task_type": "<Task selected in STEP 2>",
 "question": "<A visually grounded question in {language}>",
 "positive_answer": "<Concise, correct answer in {language}>",
 "hard_negative_answer": "<A plausible but deceptive incorrect answer in {language}>"
}

Hard Constraints:
- "task_type" must be exactly chosen from the list in STEP 2.
- Ensure the question is directly answerable from the visual or embedded textual content.
- Output ONLY the JSON object.
\end{verbatim}
\end{promptblock}

\begin{promptblock}{Video Classification Prompt}
\begin{verbatim}
You are given a video. Your job is to create ONE high-quality multimodal training example for a video classification dataset.  
The final answer MUST be a single JSON object and nothing else.

STEP 1 - Visual Analysis (less or equal than 300 {language} words)
- General overview of the video content.
- Identify primary actions, environmental settings, and the overall event type.
- Brainstorm potential ways this video could support the classification tasks listed in STEP 2.

STEP 2 - Task Selection
Choose ONE task type from the list below that best fits the video:
- Activity Recognition: Identifying the main activity or action being performed.
- Scene Parsing: Determining the primary environment or setting of the video.
- Event Categorization: Classifying the video into a specific event type or intended purpose.
- Sentiment/Intent Analysis: Recognizing the dominant emotional tone or the sentiment expressed.

STEP 3 - Populate the Example
Fill every key below using double quotes. Do not add extra keys.
{
  "description": "<STEP 1 output>",
  "task_type": "<Task Selected in STEP 2>",
  "label": "<Correct label in {language}>",
  "misleading_label": "<Plausible but incorrect label in {language} for hard negative mining>"
}

Hard Constraints:
- "task_type" must be exactly chosen from the list in STEP 2.
- "description", "label", and "misleading_label" must be in {language}.
- Output ONLY the JSON object—no extra text or explanations.
\end{verbatim}
\end{promptblock}
